
\documentclass{article}
\usepackage{colortbl}
\usepackage{makecell}
\usepackage{multirow}
\usepackage{supertabular}

\begin{document}

\newcounter{utterance}

\centering \large Interaction Transcript for game `bbh', experiment `logical\_deduction\_seven\_objects', episode 25 with reflexion{-}llama{-}3.1{-}8b{-}instruct.
\vspace{24pt}

{ \footnotesize  \setcounter{utterance}{1}
\setlength{\tabcolsep}{0pt}
\begin{supertabular}{c@{$\;$}|p{.15\linewidth}@{}p{.15\linewidth}p{.15\linewidth}p{.15\linewidth}p{.15\linewidth}p{.15\linewidth}}
    \# & \multicolumn{2}{c}{Player} && \multicolumn{2}{c}{Game Master} \\
    \hline

    \theutterance \stepcounter{utterance}  
    & & & \multicolumn{4}{p{0.6\linewidth}}{
        \cellcolor[rgb]{0.9,0.9,0.9}{
            \makecell[{{p{\linewidth}}}]{
                \texttt{\tiny{[P1$\langle$GM]}}
                \texttt{You have attempted this same task before and did not succeed. These are the plans you wrote after those attempts.} \\
\texttt{Trial \#1: New plan: } \\
\\ 
\texttt{In my previous attempt, I successfully deduced the order of the birds by breaking down the information given and using a step{-}by{-}step approach. However, I failed to consider the position of the owl in relation to the other birds. Specifically, I didn't use the fact that the blue jay is to the right of the owl (statement 5) and the raven is to the left of the falcon (statement 6) to determine the correct order.} \\
\\ 
\texttt{Next time, I will pay closer attention to the relationships between the birds and use all the given information to deduce the correct order. I will also make sure to consider the position of the owl in relation to the other birds. For example, I will use the fact that the blue jay is to the right of the owl to determine the correct order of the birds before the blue jay.} \\
\\ 
\\ 
\texttt{A logical deduction task which requires deducing the order of a sequence of objects.} \\
\\ 
\texttt{Q: The following paragraphs each describe a set of three objects arranged in a fixed order. The statements are logically consistent within each paragraph. In a golf tournament, there were three golfers: Amy, Eli, and Eve. Eve finished above Amy. Eli finished below Amy.} \\
\texttt{Options:} \\
\texttt{(A) Amy finished last} \\
\texttt{(B) Eli finished last} \\
\texttt{(C) Eve finished last} \\
\texttt{A: (B)} \\
\\ 
\texttt{Q: The following paragraphs each describe a set of three objects arranged in a fixed order. The statements are logically consistent within each paragraph. On a shelf, there are three books: a white book, a green book, and an orange book. The green book is to the right of the white book. The orange book is the rightmost.} \\
\texttt{Options:} \\
\texttt{(A) The white book is the leftmost} \\
\texttt{(B) The green book is the leftmost} \\
\texttt{(C) The orange book is the leftmost} \\
\texttt{A: (A)} \\
\\ 
\texttt{Q: The following paragraphs each describe a set of three objects arranged in a fixed order. The statements are logically consistent within each paragraph. On a shelf, there are three books: a red book, a gray book, and a white book. The white book is to the left of the gray book. The red book is the second from the left.} \\
\texttt{Options:} \\
\texttt{(A) The red book is the leftmost} \\
\texttt{(B) The gray book is the leftmost} \\
\texttt{(C) The white book is the leftmost} \\
\texttt{A: (C)} \\
\\ 
\texttt{Q: The following paragraphs each describe a set of seven objects arranged in a fixed order. The statements are logically consistent within each paragraph. On a branch, there are seven birds: a blue jay, an owl, a falcon, a hawk, a raven, a crow, and a hummingbird. The falcon is the second from the left. The crow is the rightmost. The hummingbird is to the right of the blue jay. The hawk is to the right of the hummingbird. The blue jay is to the right of the owl. The raven is to the left of the falcon.} \\
\texttt{Options:} \\
\texttt{(A) The blue jay is the fourth from the left} \\
\texttt{(B) The owl is the fourth from the left} \\
\texttt{(C) The falcon is the fourth from the left} \\
\texttt{(D) The hawk is the fourth from the left} \\
\texttt{(E) The raven is the fourth from the left} \\
\texttt{(F) The crow is the fourth from the left} \\
\texttt{(G) The hummingbird is the fourth from the left} \\
\texttt{A:} \\
            }
        }
    }
     \\ \\

    \theutterance \stepcounter{utterance}  
    & \multicolumn{4}{p{0.6\linewidth}}{
        \cellcolor[rgb]{0.561,0.737,0.561}{
            \makecell[{{p{\linewidth}}}]{
                \texttt{\tiny{[P1$\rangle$GM]}}
                \texttt{Based on the previous attempts and the given plans, I will attempt to deduce the correct order of the birds.} \\
\\ 
\texttt{From the given information:} \\
\\ 
\texttt{1. The falcon is the second from the left.} \\
\texttt{2. The crow is the rightmost.} \\
\texttt{3. The hummingbird is to the right of the blue jay.} \\
\texttt{4. The hawk is to the right of the hummingbird.} \\
\texttt{5. The blue jay is to the right of the owl.} \\
\texttt{6. The raven is to the left of the falcon.} \\
\\ 
\texttt{Using the fact that the blue jay is to the right of the owl (statement 5) and the raven is to the left of the falcon (statement 6), I will deduce the correct order of the birds before the blue jay.} \\
\\ 
\texttt{Since the raven is to the left of the falcon, and the falcon is the second from the left, the raven must be the leftmost bird.} \\
\\ 
\texttt{The order of the birds before the blue jay is: raven, falcon, owl, blue jay.} \\
\\ 
\texttt{Now, I will use the fact that the hummingbird is to the right of the blue jay (statement 3) and the hawk is to the right of the hummingbird (statement 4) to determine the correct order of the birds after the blue jay.} \\
\\ 
\texttt{The order of the birds after the blue jay is: hummingbird, hawk, crow.} \\
\\ 
\texttt{Combining the two orders, I get: raven, falcon, owl, blue jay, hummingbird, hawk, crow.} \\
\\ 
\texttt{Now, I will check the options to see which one matches the correct order.} \\
\\ 
\texttt{The correct order is: raven, falcon, owl, blue jay, hummingbird, hawk, crow.} \\
\\ 
\texttt{Comparing this order with the options, I see that none of the options match the correct order. However, I can try to find the correct answer by looking for the bird that is in the correct position in the order.} \\
\\ 
\texttt{The correct position of the blue jay is the fourth from the left. Therefore, the correct answer is:} \\
\\ 
\texttt{(A) The blue jay is the fourth from the left} \\
            }
        }
    }
    & & \\ \\

    \theutterance \stepcounter{utterance}  
    & & & \multicolumn{2}{p{0.3\linewidth}}{
        \cellcolor[rgb]{0.9,0.9,0.9}{
            \makecell[{{p{\linewidth}}}]{
                \texttt{\tiny{[GM$|$GM]}}
                \texttt{Target: (A)} \\
            }
        }
    }
    & & \\ \\

    \theutterance \stepcounter{utterance}  
    & & & \multicolumn{2}{p{0.3\linewidth}}{
        \cellcolor[rgb]{0.9,0.9,0.9}{
            \makecell[{{p{\linewidth}}}]{
                \texttt{\tiny{[GM$|$GM]}}
                \texttt{game\_result = LOSE} \\
            }
        }
    }
    & & \\ \\

    \theutterance \stepcounter{utterance}  
    & & \multicolumn{4}{p{0.6\linewidth}}{
        \cellcolor[rgb]{0.95,0.95,0.95}{
            \makecell[{{p{\linewidth}}}]{
                \texttt{\tiny{[P1$\leftrightarrow$P]}}
                \texttt{Post{-}episode Reflexion (not sent to GM):} \\
\texttt{New plan: } \\
\\ 
\texttt{In my previous attempt, I successfully deduced the order of the birds before the blue jay, but failed to consider the position of the hawk and the crow in relation to the other birds. Specifically, I didn't use the fact that the hummingbird is to the right of the blue jay (statement 3) and the hawk is to the right of the hummingbird (statement 4) to determine the correct order after the blue jay.} \\
\\ 
\texttt{Next time, I will pay closer attention to the relationships between the birds and use all the given information to deduce the correct order. I will also make sure to consider the position of the hawk and the crow in relation to the other birds. For example, I will use the fact that the hummingbird is to the right of the blue jay and the hawk is to the right of the hummingbird to determine the correct order of the birds after the blue jay.} \\
\\ 
\texttt{Additionally, I will re{-}examine the options and look for a more direct way to determine the correct answer, rather than trying to find the correct position of a single bird.} \\
            }
        }
    }
    & \\ \\

\end{supertabular}
}

\end{document}
